\documentclass[11pt]{article}
\usepackage[a4paper,margin=23mm]{geometry}
\usepackage{amsmath,amssymb,amsthm,mathtools,bm,booktabs,microtype,xurl,hyperref}
\hypersetup{colorlinks=true,linkcolor=blue,urlcolor=blue,citecolor=blue,
pdftitle={P54 scalar, common tadpole, Wilson and finite seesaw matching}}
\setlength{\emergencystretch}{2em}
\newcommand{\tr}{\operatorname{tr}}
\newcommand{\diag}{\operatorname{diag}}
\newcommand{\MS}{\overline{\mathrm{MS}}}
\newcommand{\dd}{\mathrm d}
\newcommand{\LOOP}{\mathcal L}
\newtheorem{proposition}{Proposition}
\title{Route-F P54: Scalar Matching, a Common Tadpole Anchor,\\
Factorizable Wilson Terms and Finite Seesaw Thresholds}
\author{P54PQ-v2 four-dimensional mainline}
\date{11 September 2026}
\begin{document}
\maketitle
\begin{abstract}
We calculate the upper scalar hard two-point function on the existing
Pati--Salam saddle, propagate the already fixed bosonic counterterms to
that saddle, and evaluate a factorizable scalar-exchange Wilson response.
We also calculate an exact constant-background scalar determinant
subtraction at the lower broken background and implement the universal
gauge/quartic finite seesaw contribution at moving mass thresholds.
The principal new warning is structural: the computed heavy mass
insertion has relative spectral radius $6.76$, although some projected
source corrections are much smaller. Its propagator expansion is not
controlled at the displayed benchmark. This is not a full-model
exclusion, because several matching contributions remain absent.
All results below have explicit subset boundaries; complete finite
matching, the fermion-inclusive light state and a physical fit remain open.
\end{abstract}

\section{Fixed theory, prescription and deliverable boundary}
We retain the existing renormalizable P54PQ-v2 action with
$3\,16_F+54_{H,\mathbb R}+126_{H,\mathbb C}+10_{H,\mathbb C}
+1_{H,\mathbb C}$, its two complex symmetric UV family matrices, and its
previous scalar parameter point. We neither fit a new mass formula nor
add lattice scans, geometric particle labels, or family parameters.
The upper tree stationary PS saddle is
\begin{equation}
 (w_U,\sigma_U,v_{s,U})=(1.00081030592151,0,0.254277697083708).
 \label{eq:upper}
\end{equation}
The historical broken tree vacuum is $(1,0.1265,0.25)$. All masses and
fields in numerical tables use its $\omega$ as reference. We use DR with
$d=4-2\epsilon$, $\MS$, background-field Landau gauge and the \emph{same}
previous finite fixed-VEV bosonic counterterms. DR here means dimensional
regularization, not dimensional reduction. Define $\LOOP=16\pi^2$.

The common counterterm reference is
\begin{equation}
 \mu_0=0.5626028899853225\,\omega,
 \qquad M_V^2=0.3170351789677566\,\omega^2.
\end{equation}
All new upper terms are evaluated at $\mu_0$, not the nearby historical
$\mu_U=0.5630587704385367\,\omega$. One must evolve or convert the
counterterms when changing their anchor. The upper field census is
55 positive physical heavy real scalars, 248 retained real PS scalar
coordinates, 24 gauge Goldstones and one physical PQ direction.
Retained curvatures need not be positive at this intermediate saddle.
They are not integrated as unstable massive particles.

The existing P-MATCH-KJC1 note proves the first-order transport of
kernel, sources and contact terms. This note supplies further diagrams
and tests; it does not turn transport identities into a diagram-completeness
theorem. The finite matching method is based on equality of the full and
EFT amplitudes after their common IR contributions are subtracted
\cite{henning,scalar}.

\section{Actual upper scalar hard two-point function}
\subsection{Action derivatives, not interpolation of a spectrum}
Let $X$ be the 328-component canonical real scalar coordinate, $V(X)$ the
fixed quartic potential and $\mathsf H(X)=V''(X)$. At $X=x_U$ define
\begin{equation}
 T_a=\partial_a\mathsf H(x_U),\qquad
 U_{ab}=\partial_a\partial_b\mathsf H(x_U).
\end{equation}
Because $\mathsf H$ is quadratic in $X$, for every real direction $q$
and every nonzero step $s$ the polynomial identities are
\begin{align}
 T_q&=\frac{\mathsf H(x_U+s q)-\mathsf H(x_U-s q)}{2s},\\
 U_{qq}&=\frac{\mathsf H(x_U+s q)+\mathsf H(x_U-s q)
                       -2\mathsf H(x_U)}{s^2}.
 \label{eq:jets}
\end{align}
Thus these evaluations extract the action's exact cubic and quartic
vertices, up to floating-point roundoff. They are not an extrapolated
mass fit. Steps $s=0.02$ and $0.01$ supply an independent polynomial
identity test. Hessians are evaluated by differentiating the actual
potential and cached by its source, parameters and field coordinate.

\subsection{Loop kernels and the hard-region scope}
Use a complete orthonormal scalar basis with 55 heavy mass squares $m_i^2>0$.
For the other 273 states set the propagator mass to zero \emph{only in the
leading retained-mass hard expansion}. The latter include the upper
Landau Goldstones and PQ direction. They have not been removed from the
state sum or declared physically massless by this approximation.
For $x,y\geq0$, at least one positive, define
\begin{align}
 A(x)&=x\left(\log\frac{x}{\mu_0^2}-1\right),\\
 b(x,y)&=1-\frac{x\log(x/\mu_0^2)-y\log(y/\mu_0^2)}{x-y},\\
 S(x,y)&=\int_0^1\frac{t(1-t)\,\dd t}{(1-t)x+ty}.
 \label{eq:scalar-kernels}
\end{align}
The continuous limits are
\begin{equation}
 b(x,x)=-\log(x/\mu_0^2),\quad b(x,0)=1-\log(x/\mu_0^2),
 \quad S(x,x)=\frac1{6x},\quad S(x,0)=\frac1{2x}.
\end{equation}
Expand $\tfrac12\tr\log[-\partial^2+\mathsf H(X)]$ to second order in
external scalars. The seagull comes from the term linear in $U$ and the
bubble from the term quadratic in $T$. With $\sum'$ denoting pairs with
at least one upper heavy line, the renormalized hard coefficients are
\begin{align}
 t_a^{S}&=\frac1{2\LOOP}\sum_{i\in h} A(m_i^2)T_{a,ii},\\
 \Pi_{ab}^{S}&=\frac1{2\LOOP}\left[
   \sum_{i\in h}A(m_i^2)U_{ab,ii}
   -\sum_{ij}'T_{a,ij}T_{b,ji}b(m_i^2,m_j^2)\right],\\
 k_{ab}^{S}&=\frac1{2\LOOP}\sum_{ij}'
 T_{a,ij}T_{b,ji}S(m_i^2,m_j^2).
 \label{eq:scalar-two-point}
\end{align}
These signs follow directly from the second variation of the logarithm;
the momentum derivative of the negative Euclidean bubble is positive.
Pure-soft contributions are scaleless in this hard expansion and belong
to the common EFT subtraction. No retained-mass error bound follows from
dropping them at this order.

For any real $q$, $q^Tk^Sq$ is a sum of nonnegative squares with weights
$S>0$. Hence $k^S\succeq0$ is an independent analytical check, not a
numerical assumption. The mass insertion has no such positivity theorem.

\subsection{Reconstructing all 248 real directions by symmetry}
The retained real representations are $H_8,F_{120},L_{60},R_{60}$.
Each of $H$ and $F$ is a complex field in a real-type PS representation,
and therefore consists of two equivalent real copies. A real symmetric
PS-invariant bilinear has three coefficients on each such isotypic
component. The complex-type irreducible $L$ and $R$ each contribute one.
Thus the invariant space has dimension $3+3+1+1=8$.

Construct an orthonormal basis $B_A$ of this commutant using the existing
actual conjugate-field intertwiners, not arbitrary matrices fitted to
eigenvalues. Choose eight directions for which
$P_{iA}=q_i^TB_Aq_i$ is invertible. For a calculated quadratic form $F$,
\begin{equation}
 f_A=(P^{-1})_{Ai}F(q_i,q_i),\qquad F=\sum_A f_AB_A.
\end{equation}
Two independent mixed-parent/complex-phase directions and all PS
generator commutators verify the reconstruction. The scalar kinetic
eigenvalues span
\begin{equation}
 9.38883\times10^{-5}\ \leq\lambda(k^S)\leq\ 0.00450338,
\end{equation}
whereas $\Pi^S$ spans $[-0.1408203,-0.0117783]\omega^2$.
Scalar-leg canonicalization contributes
$\delta Y^a=-\tfrac12 k^S_{ab}Y^b$, including the allowed conjugate $H/F$
invariants. It is a determined correction to the two UV spurions, not
two extra free family matrices. It must not also be included as an
external heavy leg when the same kernel is subsequently inverted.

\section{One tadpole prescription across both backgrounds}
\subsection{Why the upper tadpoles cannot be reset independently}
The previous broken-vacuum calculation fixed the bosonic finite mass
counterterms to
\begin{equation}
 (\delta\mu^2,\delta\nu^2,\delta\mu_s^2)/\omega^2
 =(-0.05633071,-5.46974409,-0.03423771).
\end{equation}
In canonical real coordinates their Hessian is
\begin{equation}
 \mathsf H_{\rm CT}=diag\left(
 -\delta\mu^2 I_{54},-\frac{\delta\nu^2}{2}I_{252},
 0_{20},-\delta\mu_s^2 I_2\right).
 \label{eq:ct}
\end{equation}
The zero $10_H$ entry does not assert that the final single-light-state
retuning is zero; that further operation has not been included here.
At $\sigma_U=0$ its tadpole cannot independently determine a second
$\delta\nu^2$. Setting both backgrounds' tadpoles to zero by independently
changing the same counterterms would silently change the prescription.

Let $R$ contain the two radial field directions $\partial X/\partial w$
and $\partial X/\partial v_s$, which span the PS-invariant heavy tadpoles.
The hard vector contribution follows by differentiating
$3\tr[M_V^4(\log(M_V^2/\mu_0^2)-5/6)]/(64\pi^2)$:
\begin{equation}
 t^V_r=\frac{3M_V^2}{2\LOOP}
 \left(\log\frac{M_V^2}{\mu_0^2}-\frac13\right)
 \tr (M_V^2)_{,r}.
\end{equation}
For the common hard-plus-counterterm radial insertion,
\begin{equation}
 t_r=t^S_r+t^V_r+(R^T\mathsf H_{\rm CT}x_U)_r,
 \qquad \delta r=-(R^T\mathsf H R)^{-1}t.
 \label{eq:radial}
\end{equation}
\begin{samepage}
The actual values are
\begin{center}
\begin{tabular}{lrr}
\toprule
Contribution & $w$ direction & $v_s$ direction\\
\midrule
Hard scalar tadpole & $-0.0220062$ & $-0.00210402$\\
Hard vector tadpole & $-0.0151893$ & $0$\\
Common counterterm & $0.1353033$ & $0.00870589$\\
Residual $t$ & $0.0981077$ & $0.00660187$\\
Linear response $\delta r$ & $-0.0139319$ & $-0.0555909$\\
\bottomrule
\end{tabular}
\end{center}
\end{samepage}
Tadpoles and responses carry $\omega^3$ and $\omega$, respectively.
The $v_s$ response is about $22\%$ of $v_{s,U}$ and is not automatically
a negligible background correction. Equation~\eqref{eq:radial} contains
only the hard-plus-CT part: retained EFT-loop tadpoles and fermionic
fixed-VEV counterterms are still absent. It is \emph{not} a completed
upper quantum stationary solution.

\subsection{Tadpole insertion from the tree heavy valley}
Let $y$ denote retained coordinates, $z$ massive coordinates and write
the tree solution as $z=z_{\rm cl}(y)$. At the PS point the mixed tree
Hessian between these inequivalent irreducible components vanishes, so
$z_{,a}=0$. Differentiating $V_{,z}(y,z_{\rm cl})=0$ twice yields
\begin{equation}
 z_{,ab}=-C^{-1}V_{,zab},\qquad C=V_{,zz}>0.
\end{equation}
Apply the chain rule to $V_{1,h}(y,z_{\rm cl})+V_{\rm CT}(y,z_{\rm cl})$:
\begin{align}
 \Delta D_{ab}&=\Pi_{ab}^{S}+\Pi_{ab}^{V}
       +(\mathsf H_{\rm CT})_{ab}+t_z^T z_{,ab},\\
 t_z^Tz_{,ab}&=\delta z^TV_{,zab},\qquad \delta z=-C^{-1}t_z.
 \label{eq:valley-mass}
\end{align}
This derives the factorizable tadpole mass term without postulating an
independently retuned upper vacuum. The two radial directions suffice
for $t_z$ by unbroken PS invariance. For the heavy six-mode propagator
below, apply the same singlet-valley differentiation while holding those
six modes as external coordinates; do not eliminate an external field
twice.
For active fields the linear vector-mass derivative vanishes, leaving
\begin{equation}
 \Pi^V=\frac{3g_U^2M_V^2}{\LOOP}
 \left(\log\frac{M_V^2}{\mu_0^2}-\frac13\right)C_S.
\end{equation}
Combining the scalar and vector kinetic terms, the canonical mass
increment is
\begin{equation}
 \Delta D_{\rm can}=\Delta D-\frac12\{k^S+k^V,D_0\}.
\end{equation}
It is real symmetric; positivity of the kinetic metric is checked.
The current canonical mass increment ranges approximately from
$-0.0260$ to $2.596\,\omega^2$. A small kinetic correction alone therefore
does not justify perturbative control of the mass or the vacuum.

\section{Factorizable upper scalar-exchange Wilson response}
\subsection{All Yukawa-coupled heavy scalar copies}
Of the 55 positive upper scalars, the actual two UV Yukawa tensors
couple to 24 real modes: four real six-dimensional PS copies. The
symmetric invariant mass block has dimension
$\dim\operatorname{Sym}^2\mathbb R^4=10$. The implementation solves
\begin{equation}
 R_i^A X_{ij}-X_{ij}R_j^A=0
\end{equation}
for the intertwiners between the six-dimensional copies using all PS
generators. Ten quadratic probes reconstruct their full hard scalar
mass block, including off-diagonal copy mixings; a separate direction
checks it. Internal scalar sums still run over all 328 states as in
Eq.~\eqref{eq:scalar-two-point}. Vector curvature, the common bosonic
counterterms and the hard valley response are then added.

Choose real Hermitian current slots $I$ from the actual fermion-pair
vertices and use
\begin{equation}
 S_2=\tfrac12\chi^TD\chi-\chi^TJI-\tfrac12 I^TCI.
\end{equation}
Our 18 slots are the real and imaginary parts of nine pairs, including
different-colour diquarks. They are probes of the scalar response, not
a complete Lorentz/gauge four-fermion operator basis. Gaussian
elimination gives the scalar-exchange contact kernel
\begin{align}
 C_0^{\rm exch}&=J_0^TD_0^{-1}J_0,\\
 \Delta C^{\rm exch}&=\Delta J^TD_0^{-1}J_0
   +J_0^TD_0^{-1}\Delta J
   -J_0^TD_0^{-1}\Delta D D_0^{-1}J_0.
 \label{eq:Wilson}
\end{align}
Here $\Delta J$ includes the previously computed pure-Yukawa and upper
vector proper vertices and fermion external legs. Heavy scalar
canonicalization is not also inserted into $J$: its compensation is
already in $D$. Explicitly, under an infinitesimal real change of heavy
field coordinates,
\begin{equation}
 \Delta D\mapsto\Delta D+s^TD_0+D_0s,\qquad
 \Delta J\mapsto\Delta J+s^TJ_0,
\end{equation}
the two added vertex terms in Eq.~\eqref{eq:Wilson} cancel the two added
propagator terms exactly. Finite differences of the full resolvent
verify Eq.~\eqref{eq:Wilson} independently.

\subsection{A non-small insertion hidden by projected channels}
For synthetic one-family choices
$(h_{\rm raw},f_{\rm raw})=(0.20+0.04i,0.025-0.01i)$ and
$(0.13-0.07i,0.041+0.013i)$, the ratios
$\|\Delta C^{\rm exch}\|_F/\|C_0^{\rm exch}\|_F$ are approximately
$0.0441$ and $0.361$. Neither is a fitted amplitude or proton lifetime.
More importantly, define
\begin{equation}
 E=D_0^{-1/2}\Delta D D_0^{-1/2}.
\end{equation}
The actual computed subset has
\begin{equation}
 \lambda(E)\in[-0.1314683,6.7559262],\qquad \rho(E)=6.7559262.
 \label{eq:large-insertion}
\end{equation}
On its dominant unit eigenvector $u$, the contractions
$u^TD_0^{-1/2}\Delta D_XD_0^{-1/2}u$ from scalar loops, vector loops,
the common bosonic counterterm and the valley response are respectively
$(-0.538946,-0.016987,+7.346374,-0.034515)$. These add to $6.755926$;
they are not eigenvalues of the separate contributions. The large
direction is therefore dominated by the transported finite counterterm,
not a loss of numerical precision in the loop kernels.
\begin{proposition}[Exact convergence criterion for this subset]
For $D_0>0$ and real symmetric $\Delta D$, the Taylor/Neumann series of
$(D_0+\ell\Delta D)^{-1}$ about $\ell=0$ converges in matrix norm if
$|\ell|\rho(E)<1$, and fails when $|\ell|\rho(E)>1$.
\end{proposition}
\begin{proof}
Factor the inverse as $D_0^{-1/2}(I+\ell E)^{-1}D_0^{-1/2}$ and
orthogonally diagonalize $E$. Each eigenvalue gives an ordinary geometric
series with ratio $-\ell\lambda(E)$. The largest absolute eigenvalue
sets the radius of convergence.
\end{proof}
The radius is about $0.1480$, so the displayed subset is not a small
insertion at $\ell=1$. Positivity of $I+E$ does not repair convergence;
both properties can and here do coexist. Small source projections can
miss the large-eigenvalue directions and cannot certify the full kernel.

This is a precise obstruction to promoting the present fixed-parameter
subset, not a proof of nonperturbativity or exclusion of the complete
theory. Missing finite diagrams, fermionic counterterms, parameter
feedback and retained-mass corrections may modify $E$. Resumming only
this subset is not a new physical prediction. Nonfactorizable boxes,
direct vector exchange, scalar three/four-point Wilson vertices and a
complete operator basis remain uncomputed.

\section{Lower constant-background finite scalar subtraction}
At the actual broken background, let $B_A$ span the 248 upper retained
coordinates and $B_H$ the 55 upper heavy coordinates. Write
\begin{equation}
 A=B_A^T\mathsf HB_A,\quad C=B_H^T\mathsf HB_H,\quad
 W=B_H^T\mathsf HB_A,\quad S_0=A-W^TC^{-1}W.
\end{equation}
With $O$ the 24 upper gauge-orbit directions, the horizontal projector
is $P=I-O(O^TO)^{-1}O^T$. At this background $B_H^TO=0$, so in the
303-dimensional chart $(B_A,B_H)$ the scalar kinetic metric is
$\diag(G_0,I_{55})$ with $G_0=B_A^TPB_A$. The PQ spectator is not part
of this chart. The exact scalar Schur kernel at Euclidean $z=p_E^2$ is
\begin{align}
 S(z)&=A+zG_0-W^T(C+zI)^{-1}W,\\
 S_{[2]}(z)&=S_0+zG,\qquad G=G_0+W^TC^{-2}W.
 \label{eq:nonlocal}
\end{align}
Block Gaussian elimination gives the pointwise identity
\begin{equation}
 \log\det(\mathsf M+z\mathsf G)
 =\log\det(C+zI)+\log\det S(z).
 \label{eq:det}
\end{equation}
The full generalized scalar spectrum has 290 positive modes and 13
zeros; the local 248-dimensional operator has 235 positive modes and
13 zeros. The zeros include nine lower gauge directions and four
light-Higgs real directions. They are not thirteen physical massless
particles. The 290 positive masses agree with the original real Hessian.

Let
\begin{equation}
 v(x;\mu_0)=\frac{x^2}{64\pi^2}
       \left(\log\frac{x}{\mu_0^2}-\frac32\right),\qquad v(0)=0.
\end{equation}
Denote the generalized full, upper-$C$, and local spectra by $m_f^2$,
$m_C^2$ and $m_l^2$. In dimensional regularization, field-independent
kinetic normalization factors multiply scaleless integrals and drop out.
Integrating Eq.~\eqref{eq:det} in this scalar sector yields
\begin{align}
 V^{(1)}_{\rm PS,nonlocal}&=\sum_fv(m_f^2)-\sum_Cv(m_C^2),\\
 \Delta V^{(1)}_{\rm loc}&=
 \sum_fv(m_f^2)-\sum_Cv(m_C^2)-\sum_lv(m_l^2).
 \label{eq:finite-subtraction}
\end{align}
Numerically, in $\omega^4$ units,
\begin{align}
 V_{\rm full}&=-0.04307001295,& V_C&=-0.00596443931,\\
 V_{\rm PS,nonlocal}&=-0.03710557365,&
 V_{\rm PS,local}&=-0.03705904935,\\
 \Delta V^{(1)}_{\rm loc}&=-4.65242986\times10^{-5}.&&
\end{align}
Six momentum checks verify Eq.~\eqref{eq:det}; a transformed broken PS
background gives the same finite difference. Its scale derivative is
independently fixed by
\begin{equation}
 \frac{\dd\Delta V^{(1)}_{\rm loc}}{\dd\log\mu}
 =-\frac{\sum_fm_f^4-\sum_Cm_C^4-\sum_lm_l^4}{32\pi^2}.
\end{equation}

The conclusion is not that this vacuum constant is an observable or a
complete lower threshold. It is that replacing the exact nonlocal
operator by its two-derivative expansion \emph{inside an unrestricted
loop integral} loses a definite finite term. Higher-derivative/tree
Wilson insertions and the associated subtraction must be tracked.
Background/source derivatives, transverse gauge fields, ghosts,
Goldstone gauge-fixing terms and finite Yukawa/box contributions remain
necessary for the full gauge-covariant matching functional. The quotient
calculation alone is not that functional.

\section{Finite seesaw: complete restricted limit and actual subset}
\subsection{Conventions and the terminal canonical limit}
We use $Y_{\alpha i}$ with lepton rows, $M_R^T=M_R$, and
\begin{equation}
 U^TM_RU=\diag(M_i),\quad M_i>0,\qquad
 Q=YM^{-1}Y^T,\qquad C_{5,\rm tree}^{I}=-Q.
\end{equation}
Our $C_5$ sign is opposite to that of Ref.~\cite{zhangzhou}; our $g_1$
is GUT normalized, $g_Y^2=3g_1^2/5$. Write
$\ell_i=\log(\mu^2/M_i^2)$. The canonical type-I one-loop result can be
organized into the following finite contributions
\cite{zhangzhou}:
\begin{align}
 z_H&=\frac1{\LOOP}\sum_i\|Y_i\|^2(\tfrac12+\ell_i),\\
 k_{L,\rm row}&=\frac1{\LOOP}Y\diag\left(\frac{3+2\ell_i}{4}\right)Y^\dagger,\\
 \delta C_g&=-\frac1{\LOOP}\sum_i
 \left[2\lambda(1+\ell_i)+\frac{g_Y^2+g_2^2}{4}(1+3\ell_i)\right]
 \frac{Y_iY_i^T}{M_i},\label{eq:finiteC}\\
 \delta C_{\rm legs}&=-z_H C_{5,\rm tree}^{I}
 -\frac12\left(k_{L,\rm row}C_{5,\rm tree}^{I}
              +C_{5,\rm tree}^{I}k_{L,\rm row}^T\right).
\end{align}
The row tensor is the transpose of the literal left-Weyl kinetic matrix;
this distinction matters under complex family transformations. The
terminal canonical result is $C_5=-Q+\delta C_g+\delta C_{\rm legs}$.
Its applicability here is restricted to all sterile fields removed,
$C_{HN}=0$, no pre-existing $C_5$, and leading $q_H/M_i^2=0$.
It is a complete $C_5$ formula in that limit, not complete SMEFT matching.
At $\mu=M$ the independent degenerate wave-function calculation gives
the factors $1/2$ and $3/4$ \cite{ohlsson}.

\subsection{Independent RG and covariance checks}
At fixed parameters let $A_Y=YY^\dagger$. Differentiating the finite
terms gives
\begin{align}
 \frac{\dd\delta C_g}{\dd\log\mu}
 &=-\frac{4\lambda+\tfrac32(g_Y^2+g_2^2)}{\LOOP}\,Q,\label{eq:RGg}\\
 \frac{\dd\delta C_{\rm legs}}{\dd\log\mu}
 &=\frac{2\tr(A_Y)Q+\tfrac12(A_YQ+QA_Y^T)}{\LOOP}.
 \label{eq:RGchecks}
\end{align}
These expressions equal the independently implemented
$\beta_{C_5}^{\rm below}-\beta_{(C_5-Q)}^{\rm above}$ in the canonical
limit. In the active theory,
\begin{equation}
 \beta_{C_5-Q}=\beta_{C_5}-\beta_YM^{-1}Y^T-YM^{-1}\beta_Y^T
                    +YM^{-1}\beta_M M^{-1}Y^T.
\end{equation}
Keeping the matrix ordering is essential. Finite differences in
$\log\mu$ test the relation without differentiating eigenvectors.
Under $Y\mapsto U_L^TYU_N$, $M\mapsto U_N^TMU_N$, the result transforms
as $C_5\mapsto U_L^TC_5U_L$. Exactly degenerate positive Takagi blocks
permit a real orthogonal residual transformation, under which every
block sum is invariant. These tests verify generic complex families
and an independent $O(3)$ block rotation.

\subsection{What can be applied to the actual P54 trajectory}
The actual P54 scalar elimination generates $C_{HN}\ne0$ and a type-II
Weinberg term. We explicitly reject applying the restricted terminal
formula to that input. Instead, the new sequential callback adds only
the universal $\delta C_g$ of Eq.~\eqref{eq:finiteC} for the removed
Takagi block. The existing nonzero $C_{HN}$ remains in running and the
tree block-Schur pullback. No missing finite term is assigned zero and
then labelled a complete result.

For $M_R=\left(\begin{smallmatrix}M_{rr}&M_{rh}\\M_{hr}&M_{hh}\end{smallmatrix}\right)$,
the exact tree map is
\begin{align}
 M_r'&=M_{rr}-M_{rh}M_{hh}^{-1}M_{hr},\\
 Y_r'&=Y_r-Y_hM_{hh}^{-1}M_{hr},\qquad
 C_5'=C_5-Y_hM_{hh}^{-1}Y_h^T.
\end{align}
Its field-dependent derivative gives the retained $C_{HN}$ pullback.
For a moving Takagi block, $M_{rh}=0$ at the matching point but this
covariant general formula is still the implementation's regression
reference. For any removed block $h$, the finite gauge/quartic derivative
is Eq.~\eqref{eq:RGg} with $Q$ replaced by $Q_h$. It equals the
gauge/quartic part of the corresponding active-theory beta difference.

The threshold event is solved from the running matrix,
\begin{equation}
 \mu=\kappa\,s_{\max}\big(M_R(\mu)\big),
 \qquad \kappa=1\quad\hbox{for the central diagnostic}.
\end{equation}
Degenerate blocks are removed together. Event ordering is recalculated
after each evolution and finite matching step, not frozen from the
starting spectrum. Tests also vary $\kappa$ to $0.8$ and $1.25$.
The two existing synthetic families run from $0.1\omega$ to
$10^{-4}\omega$ through three live finite events. Relative terminal
$C_5$ changes versus the old tree-threshold trajectories are $0.00221120$
and $0.00224041$. These are diagnostic changes, not fitted neutrino
observables or calibrated error estimates. A finite Wilson coefficient
need not be continuous: matching equates amplitudes, not raw coefficients.

Missing finite insertions of $C_{HN}$ and pre-existing $C_5$, mixed
removed/retained sterile graphs, and correlated finite $Y_\nu,M_R$,
Higgs and wave-function maps prevent complete P54 finite seesaw closure.

\section{Verification ledger and next decisions}
\begin{center}
\begin{tabular}{lr}
\toprule
Verifier (new unless indicated) & Checks passed\\
\midrule
Upper scalar hard matching / common tadpole & 14/14\\
Upper factorizable scalar-exchange Wilson & 16/16\\
Lower finite scalar determinant subtraction & 22/22\\
Finite seesaw limit and live universal subset & 24/24\\
Original sequential-seesaw default regression & 46/46\\
\bottomrule
\end{tabular}
\end{center}
Passing these checks certifies the stated algebra and numerical tests,
\emph{not} diagram completeness, perturbativity or empirical validity.
Each new JSON report records source hashes, numerical matrices and
explicit false full-matching/fit flags. New polynomial vertices were
computed only where the original cache lacked them; no new stationary
point or lattice scan was performed.

The next work should address the structural obstruction before a fit:
\begin{enumerate}
\item Complete the common-prescription tadpole/parameter bookkeeping and
remaining upper Wilson diagrams, including fermionic counterterms when
actual family matrices are supplied. Separate source-projected smallness
from the full mass-operator test in Eq.~\eqref{eq:large-insertion}.
If this benchmark remains uncontrolled, compare a consistently converted
renormalization prescription or reject the benchmark; do not simply
resum one subset or adjust a mass formula.
\item Differentiate the lower source-dependent hard-minus-EFT functional
and include its gauge/ghost/Goldstone, Yukawa and contact diagrams. Use
the exact scalar subtraction as a regression, not as a replacement for
the missing functional. Bound the retained-mass expansion or keep the
necessary mass dependence.
\item Complete finite seesaw in the actual EFT with $C_{HN}$, type-II
$C_5$ and partially retained sterile fields, using the canonical limit
and beta-difference identities as independent tests.
\item Only after matching and perturbative/background control, solve the
same-prescription full complex light state, enforce one light-doublet
condition and the heavy-block gap, then attempt the constrained fit.
\end{enumerate}
The action remains frozen throughout this checkpoint. No physical-fit,
determinant, portal or unconditional first-principles closure is promoted.

\subsection*{Reproducible artifacts}
All paths below are relative to \texttt{route\_f/}. Each verifier has an
identically named JSON/Markdown pair in \texttt{output/}, without its
\texttt{verify\_} prefix:
\begin{itemize}
\item \path{code/verify_p54_upper_scalar_matching.py};
\item \path{code/verify_p54_upper_wilson_exchange.py};
\item \path{code/verify_p54_lower_scalar_finite_subtraction.py};
\item \path{code/verify_p54_finite_seesaw_matching.py}.
\end{itemize}
The first three accept \texttt{--cache-dir} pointing to the existing
\path{tmp/p54_full_doublet_cw} cache in the original checkout. Newly
needed polynomial Hessians are written to this worktree's
\path{tmp/p54_upper_scalar_jets}. Runtime versions are NumPy 2.0.2,
SciPy 1.13.1 and JAX/JAXlib 0.4.30. The original default seesaw regression
was rerun; its tree-threshold report was refreshed with the current source
hash while keeping its historical model/date label. It is distinct from
the new finite-threshold report.

\begin{thebibliography}{9}
\bibitem{henning} B. Henning, X. Lu and H. Murayama,
One-loop matching and running with covariant derivative expansion,
\href{https://arxiv.org/abs/1604.01019}{arXiv:1604.01019}.
\bibitem{scalar} M. Gabelmann, M. M\"uhlleitner and F. Staub,
Automatised matching between two scalar sectors at the one-loop level,
\href{https://doi.org/10.1140/epjc/s10052-019-6570-5}{Eur. Phys. J. C 79, 163 (2019)}.
\bibitem{zhangzhou} D. Zhang and S. Zhou,
Complete one-loop matching of the type-I seesaw model onto the Standard
Model effective field theory, especially Eqs.~49, 51, 89 and 91,
\href{https://arxiv.org/html/2107.12133v3}{arXiv:2107.12133v3}.
\bibitem{ohlsson} T. Ohlsson and M. Pernow,
One-loop matching conditions in neutrino effective theory,
especially Eqs.~42 and 45,
\href{https://arxiv.org/html/2201.00840v2}{arXiv:2201.00840v2}.
\end{thebibliography}
\end{document}
